\section{Monomial symmetric functions and powers sums}\label{Section3}

Fix any f\/ield $\mathbb F$.
Denote by $\Lambda$ the $\mathbb F$-algebra of symmetric functions in inf\/initely many variables
$x_1,x_2,\ldots$.
Following~\cite{M} we will introduce some standard bases of $\Lambda$.

Let $\lambda=(\lambda_1,\lambda_2,\ldots)$ be any partition of $0,1,2,\ldots$.
Here we assume that $\lambda_1\ge\lambda_2\ge\cdots$.
The number of non-zero parts is called the {\it length} of $\lambda$ and denoted by $\ell(\lambda)$.
We will also denote $|\lambda|=\lambda_1+\lambda_2+\cdots$ as usual.
Let $m_\lambda\in\Lambda$ be the \textit{monomial symmetric function} corresponding to the partition
$\lambda$.
By def\/inition,
\begin{gather}
\label{mon}
m_\lambda=\sum_{i_1<\dots<i_k}
\sum_{\sigma\in\mathfrak{S}_k}d_\lambda^{-1}x_{i_{\sigma(1)}}^{\lambda_1}\cdots x_{i_{\sigma(k)}}^{\lambda_k},
\end{gather}
where we write $k$ instead of $\ell(\lambda)$.
Here $\mathfrak{S}_k$ is the symmetric group permuting $1,\ldots, k$ while
\begin{gather}
\label{cela}
d_\lambda=k_1!k_2!\cdots
\end{gather}
if $k_1, k_2,\ldots$ are the respective multiplicites of the parts $1,2,\ldots$ of $\lambda$.
If $\ell(\lambda)=0$ then $m_\lambda=1$.
The monomial symmetric functions corresponding to partitions of $0,1,2,\ldots$ form a~basis of~$\Lambda$.

For $n=1,2,\ldots$ let $p_n\in\Lambda$ be the \textit{power sum symmetric function} of degree~$n$.
By def\/inition,
\begin{gather}
\label{pon}
p_n=x_1^n+x_2^n+\cdots.
\end{gather}
More generally, for any partition $\lambda$ put $ p_\lambda=p_{\lambda_1}\cdots p_{\lambda_k} $ where
$k=\ell(\lambda)$ as in~\eqref{mon}.
The ele\-ments~$p_\lambda$ form another basis of $\Lambda$.
In other words, the elements $p_1,p_2,\ldots$ are free generators of the commutative algebra~$\Lambda$ over
$\mathbb F$.
We will also use the formal power series in the variable $t$
\begin{gather*}
P(t)=p_1+p_2t+p_3t^2+\cdots.
\end{gather*}

Def\/ine a~bilinear form $\langle\;,\;\rangle$ on the vector space $\Lambda$ by setting for any two
partitions $\lambda$ and $\mu$
\begin{gather}
\label{schurprod}
\langle p_\lambda,p_\mu\rangle=z_\lambda\delta_{\lambda\mu},
\qquad
\text{where}
\qquad
z_\lambda=1^{k_1}k_1!2^{k_2}k_2!\cdots
\end{gather}
in the notation~\eqref{cela}.
This form is obviously symmetric and non-degenerate.
We will indicate by the superscript ${}^\perp$ the operator conjugation relative to this bilinear form.
In particular, by~\eqref{schurprod} for the operator conjugate to the multiplication in $\Lambda$ by $p_n$
with $n\ge1$ we have $ p_n^{\perp}=n\partial /\partial p_n.
{}$

The basis of $p_\lambda$ is related to the basis of monomial symmetric functions as follows.
For any two partitions $\lambda$ and $\mu$ let $r_{\lambda\mu}$ be the number of maps
$\theta:\{1,\ldots,\ell(\mu)\}\to\{1,2,\ldots\}$~such~that
\begin{gather}
\label{refine}
\sum_{\theta(j)=i}\mu_j=\lambda_i
\qquad
\text{for each}
\quad
i=1,2,\ldots.
\end{gather}
For any such $\theta$ the partition $\mu$ in~\eqref{refine} is called a~\textit{refinement} of $\lambda$.
Note that if $r_{\lambda\mu}\neq0$ then $|\lambda|=|\mu|$.
Moreover, then by [I.6.10] we have $\mu\le\lambda$ in the \textit{natural partial ordering} of partitions:
\begin{gather*}
\mu_1\le\lambda_1,\qquad \mu_1+\mu_2\le\lambda_1+\lambda_2,\qquad \ldots.
\end{gather*}
By [I.6.9] we have the relation
\begin{gather}\label{pm}
p_\mu=\sum_\mu r_{\lambda\mu}m_\lambda.
\end{gather}

\section{Elementary and complete symmetric functions}\label{Section4}

For $n=1,2,\ldots$ let $e_n\in\Lambda$ be the \textit{elementary symmetric function} of degree~$n$.
By def\/inition,
\begin{gather*}
e_n=\sum_{i_1<\dots<i_k}x_{i_1}\cdots x_{i_k}.
\end{gather*}
Note that by the def\/inition~\eqref{mon} we have $e_n=m_\lambda$ where $\lambda=(1^n)$.
Further, let $h_n\in\Lambda$ be the \textit{complete symmetric function} of degree $n$.
By def\/inition,
\begin{gather*}
h_n=\sum_{|\lambda|=n}m_\lambda.
\end{gather*}
We will also use the formal power series in the variable $t$
\begin{gather*}
E(t)=1+e_1t+e_2t^2+\cdots
\qquad
\text{and}
\qquad
H(t)=1+h_1t+h_2t^2+\cdots.
\end{gather*}
In particular, by [I.2.6] we have the relation
\begin{gather}
\label{eh}
E(-t)H(t)=1.
\end{gather}
Further, by [I.2.10] we have
\begin{gather}
\label{hep}
P(t)=E^{\prime}(-t)/E(-t).
\end{gather}

For any partition $\lambda$ denote $ e_\lambda=e_{\lambda_1}e_{\lambda_2}\cdots $ and $
h_\lambda=h_{\lambda_1}h_{\lambda_2}\cdots $ where $e_0=h_0=1$.
The elements $e_\lambda$ form a~basis of $\Lambda$, and so do the elements~$h_\lambda$.
It follows from~\eqref{eh} that for any~$n\ge1$
\begin{gather}
\label{ehnu}
(-1)^ne_n=\sum_{|\lambda|=n}(-1)^{\ell(\lambda)}h_\lambda\ell(\lambda)
!/d_\lambda.
\end{gather}
Furthermore, according to [I.4.5] the bases of $m_\lambda$ and $h_\lambda$ are dual to each other relative
to the bilinear form~\eqref{schurprod} on $\Lambda$:
\begin{gather}
\label{mh}
\langle m_\lambda,h_\mu\rangle=\delta_{\lambda\mu}.
\end{gather}
It follows from this duality that for any partitions $\lambda$ and $\nu$ we have
\begin{gather}
\label{hperp}
h_\nu^\perp(m_\lambda)=
\begin{cases}
m_\mu&\text{if}\quad\lambda=\mu\sqcup\nu,
\\
0&\text{otherwise,}
\end{cases}
\end{gather}
where $\mu\sqcup\nu$ is the partition obtained by collecting the parts of~$\mu$ and~$\nu$ together.
For~the~opera\-tor~$p_n^{\perp}$ with $n\ge1$ acting on the elements of~$\Lambda$ expressed as
polynomials in $h_1, h_2, \ldots$ by [Ex.~I.5.3]
\begin{gather}
\label{pperp}
p_n^{\perp}=\sum_{k\ge0}h_k\partial /\partial h_{n+k}.
\end{gather}

\section{Jack symmetric functions}\label{Section5}

Now let $\mathbb F$ be the f\/ield $\mathbb Q(\alpha)$ where $\alpha$ is another variable.
Generalizing~\eqref{schurprod} def\/ine a~bilinear form $\langle\;,\;\rangle_\alpha$ on $\Lambda$
by setting for any partitions $\lambda$ and $\mu$
\begin{gather}
\label{jackprod}
\langle p_\lambda,p_\mu\rangle_\alpha=\alpha^{\ell(\lambda)}
z_\lambda\delta_{\lambda\mu}.
\end{gather}
This form is symmetric and non-degenerate.
In particular, if $\alpha=1$ then this is the form $\langle\;,\; \rangle$.
We will indicate by the superscript ${}^\ast$ the operator conjugation relative to the form
$\langle\;,\;\rangle_\alpha$.
Using the def\/inition~\eqref{jackprod} for the operator conjugate to the multiplication in $\Lambda$ by
$p_n$ with $n\ge1$ we~get
\begin{gather}
\label{pnast}
p_n^{\ast}=\alpha n\partial /\partial p_n.
\end{gather}

By [Ex.~VI.4.2] there exists a~unique family of elements $J_\lambda\in\Lambda$ such that
\begin{gather}
\label{plamu}
\langle J_\lambda,J_\mu\rangle_\alpha=0
\qquad
\text{for}
\quad
\lambda\neq\mu
\end{gather}
and such that the $J_\lambda$ equals $m_\lambda$ plus a~linear combination of the elements $m_\mu$ with
$\mu<\lambda$ in the natural partial ordering.
The elements $J_\lambda\in\Lambda$ are called the \textit{Jack symmetric functions}.
In the case $\alpha=1$ the $J_\lambda$ coincides with the \emph{Schur symmetric function}~$s_\lambda$.

The Jack symmetric functions $J_\lambda$ form a~basis of the vector space $\Lambda$ over the f\/ield
$\mathbb Q(\alpha)$.
Let us now def\/ine the operators $A^{(1)},A^{(2)},\ldots$ acting on
$\Lambda$ as follows.
Put
\begin{gather}
\label{au}
A(u)=1+A^{(1)}/(u)_1+A^{(2)}/(u)_2+\cdots,
\end{gather}
where $u$ is another variable and $(u)_k=u(u+1)\cdots(u+k-1)$ is the \emph{Pochhammer symbol}.
Then
\begin{gather}
\label{ajack}
A(u)J_\lambda=\prod_{i\ge1}\frac{u+i-1-\alpha\lambda_i}{u+i-1}\cdot J_\lambda
\end{gather}
by def\/inition.
Note that in the inf\/inite product displayed above the only factors dif\/ferent from~$1$ are those
corresponding to $i=1,\ldots,\ell(\lambda)$.
In our recent publication~\cite{NS1} we proved the following theorem.
Another proof can be obtained by using the results of~\cite{Shi}, see also~\cite{NS2}.
\begin{theorem}\label{Theorem1}
In the notation~\eqref{cela} for each $k=1,2,\ldots$ we have
\begin{gather}
\label{basic}
A^{(k)}=(-1)^k\sum_{\ell(\lambda)=k}d_\lambda m_\lambda m_\lambda^{\ast},
\end{gather}
where the sum is taken over all partitions $\lambda$ of the fixed length $k$.
\end{theorem}

The Jack symmetric functions are eigenvectors of the operators
$A^{(1)},A^{(2)},\ldots$ by def\/inition.
In particular, these operators are self-conjugate relative to the form $\langle\;,\;\rangle_\alpha$ by the
orthogonality condition~\eqref{plamu}.
This self-conjugacy is also transparent from the explicit formula~\eqref{basic}.

By inverting the relation~\eqref{pm} any monomial symmetric function $m_\lambda$ can be expressed as
a~linear combination of the functions $p_\mu$ where $\lambda$, $\mu$ are partitions of the same number and~$\lambda\le\mu$.
By substituting into~\eqref{basic} and using~\eqref{pnast}, one can write each operator
$A^{(k)}$ in terms of~$p_n$ and $\partial /\partial  p_{n}$ where $n=1,2,\ldots$.
However, the result of this procedure for arbitrary $k$ does not seem to be suf\/f\/iciently explicit.
In the next section we will overcome this dif\/f\/iculty by considering another collection of operators on
$\Lambda$ such that the Jack symmetric functions are their eigenvectors.

\section{Lax matrix}\label{Section6}

Let us def\/ine the operators $I^{(1)},I^{(2)},\ldots$ acting on the vector
space $\Lambda$ over $\mathbb Q(\alpha)$ by the equations~\eqref{iu} and~\eqref{hdef}.
In particular, then we have the relations
\begin{gather*}
I^{(1)}=-A^{(1)}
\qquad
\text{and}
\qquad
I^{(2)}=A^{(1)}\big(A^{(1)}
+1\big)-2A^{(2)}.
\end{gather*}
Note that then due to the def\/inition~\eqref{ajack} of $A(u)$ we also have the equality
\begin{gather*}
I(u)J_\lambda=uJ_\lambda-(u+\ell(\lambda))\prod_{1\le i
\le\ell(\lambda)}\frac{u+i-1-\alpha\lambda_i}{u+i-\alpha\lambda_i}\cdot J_\lambda.
\end{gather*}

Let us now introduce the inf\/inite matrix
\begin{gather}
\label{Lax_matrix}
L=
\begin{bmatrix}
(\alpha-1)&p_1&p_2&p_3&\hspace{-1pt}\ldots
\\
p_1^*&2(\alpha-1)&p_1&p_2&\ddots
\\
p_2^*&p_1^*&3(\alpha-1)&p_1&\ddots
\\[-2pt]
\raisebox{-1pt}{\vdots}\hspace{12pt}\ddots\hspace{-20pt}&\hspace{24pt}\ddots\hspace{-20pt}
&\hspace{24pt}\ddots\hspace{-20pt}&\hspace{24pt}\ddots\hspace{-20pt}
\end{bmatrix},
\end{gather}
where the rows and columns can be labelled by the indices $i,j=1,2,\ldots$.
We shall call it the \emph{Lax matrix}.
It is convenient to set $p_{-n}=p^{\ast}_n$ for $n\ge1$.
We also keep assuming that $p_{0}=0$.
Then $L=[ L_{ij}]_{i,j=1}^{\infty}$ where
\begin{gather}
\label{Ljk}
L_{ij}=j(\alpha-1)\delta_{ij}+p_{j-i}.
\end{gather}
Note that the matrix $L$ is self-conjugate relative to the form $\langle\; ,\; \rangle_\alpha$.
Also introduce the inf\/inite row vector $p=[p_1\; p_2\; \ldots]$ and its conjugate column vector
\begin{gather*}
p^{\ast}=
\begin{bmatrix}
p^{\ast}_1
\\[2pt]
p^{\ast}_2
\\[6pt]
\raisebox{0pt}[0pt][2pt]{\vdots}
\end{bmatrix}
.
\end{gather*}
The main result of the present article is the next theorem which provides an explicit expression for every
operator $I^{(k)}$ in terms of $p_n$ and $\partial /\partial  p_{n}$ where
$n=1,2,\ldots$.
\begin{theorem}\label{Theorem2}
We have the equality
\begin{gather}
\label{hua}
I(u)=p(u-L)^{-1}p^{\ast},
\end{gather}
where the inverse to the infinite matrix $u-L$ is regarded as formal power series in $u^{-1}$:
\begin{gather*}
(u-L)^{-1}=u^{-1}+Lu^{-2}+L^2u^{-3}+\cdots.
\end{gather*}
\end{theorem}

Note that each coef\/f\/icient of the series in $u^{-1}$ at the right hand side of the
equality~\eqref{hua} is an inf\/inite sum of operators acting on $\Lambda$.
However, only f\/inite number of the summands do not vanish on any subspace in $\Lambda$ of a~f\/ixed
degree in $x_1,x_2,\ldots$.
This can be easily seen by restating Theorem~\ref{Theorem2} as the equality for every index $k\ge1$
\begin{gather}
\label{hk}
I^{(k)}=\sum_{i_1,\ldots,i_k=1}^\infty p_{i_1}
L_{i_1i_2}\cdots L_{i_{k-1}i_k}p^{\ast}_{i_k}
\end{gather}
and by performing summation in~\eqref{hk} consecutively over the indices $i_k,i_{k-1},\ldots,
i_1$.
In the next section we explain our method of proving Theorem~\ref{Theorem2}.
The proof will be completed after that.

\section[Baker-Akhiezer vector]{Baker--Akhiezer vector}\label{Section7}

Consider the inf\/inite column vector
\begin{gather}
\label{psiu}
\Psi(u)=
\begin{bmatrix}
\raisebox{0pt}[12pt][0pt]{$\Psi_1(u)$}
\\
\Psi_2(u)
\\
\vdots
\end{bmatrix}
=(u-L)^{-1}p^{\ast}A(u+1).
\end{gather}
Every entry $\Psi_n(u)$ of this vector is a~formal power series in $u^{-1}$ with coef\/f\/icients acting on
$\Lambda$; see the remark in the end of previous section.
We shall call it the \emph{Baker--Akhiezer vector} for the Lax matrix~\eqref{Lax_matrix}, since in the
classical limit this is an analog of the Baker--Akhiezer function for a~rational spectral curve~\cite{Skl}.
Our proof of Theorem~\ref{Theorem2} is based on the explicit formula for $\Psi_n(u)$ given by the next theorem.
This theorem is another principal result of the present article.
Let
\begin{gather*}
\Psi_n(u)=\Psi^{(1)}_n/(u+1)_1+\Psi^{(2)}_n/(u+1)_2+\cdots,
\end{gather*}
where the coef\/f\/icients $\Psi^{(1)}_n,\Psi^{(2)}_n,\ldots$ are operators
acting on the vector space $\Lambda$ over $\mathbb Q(\alpha)$.
\begin{theorem}\label{Theorem3}
For any indices $k,n\ge1$ we have the equality
\begin{gather}
\label{psink}
\Psi^{(k)}_n=(-1)^{n+k}\sum_{\ell(\lambda)=k}d_\lambda e_n^{
\perp}(m_\lambda)m_\lambda^{\ast}.
\end{gather}
\end{theorem}

Note that by using the relation~\eqref{ehnu} and then~\eqref{hperp} here we can write
\begin{gather}\label{firstfactor}
(-1)^ne_n^{\perp}(m_\lambda)
=\sum\limits_{\substack{\lambda=\mu\sqcup\nu\\|\nu|=n}}(-1)^{\ell(\nu)}m_\mu\ell(\nu)!/d_\nu,
\end{gather}
where $\nu$ ranges over all partions of $n$ such that by taking the parts of $\nu$ together with those of
another partition $\mu$ we get $\lambda$.
The partition $\mu$ is then determined by $\lambda$ and $\nu$ uniquely.

Theorem~\ref{Theorem3} shows that the series $\Psi_n(u)$ can be obtained from the series $A(u+1)$ by applying the
operator $(-1)^ne_n^{\perp}$ to the factor $m_\lambda$ in the
coef\/f\/icient~\eqref{basic}.
Thus $\Psi_n(u)$ can be expressed in terms of $A(u+1)$ without using Theorem~\ref{Theorem1}.
However we will prove Theorem~\ref{Theorem3} as stated here.

Theorem~\ref{Theorem2} will then follow rather easily.
Indeed, due to the def\/inition~\eqref{psiu} of $\Psi(u)$ to obtain the stated equality~\eqref{hua} it
suf\/f\/ices to prove that
\begin{gather*}%\label{pha}
p\Psi(u)=I(u)A(u+1).
\end{gather*}
By using the def\/inition~\eqref{hdef} of $I(u)$ the last displayed equality can be rewritten as
\begin{gather}
\label{pepsi}
p\Psi(u)=uA(u+1)-uA(u).
\end{gather}
Here on the left hand side we have a~product of a~row vector by a~column vector:
\begin{gather*}
p\Psi(u)=\sum_{n\ge1}p_n\Psi_n(u).
\end{gather*}
By using Theorems~\ref{Theorem1} and~\ref{Theorem3} along with the obvious relation for any $k\ge1$
\begin{gather*}
u/(u+1)_k-u/(u)_k=-k/(u+1)_k,
\end{gather*}
the equality~\eqref{pepsi} follows from the proposition which is
stated and proved below.
\begin{proposition}\label{Prop1}
For any partition $\lambda$ of the length $k\ge1$ we have
\begin{gather}\label{pro1}
\sum_{n\ge1}(-1)^np_ne_n^{\perp}(m_\lambda)=-km_\lambda.
\end{gather}
\end{proposition}
\begin{proof}
Using the bilinear form $\langle\;,\;\rangle$ take the operator on $\Lambda$ conjugate to the operator which
is applied to $m_\lambda$ at the left hand side of the equality~\eqref{pro1}.
The conjugate operator equals
\begin{gather*}
\sum_{n\ge1}(-1)^ne_np_n^{\perp}
=\sum_{n\ge1}\sum_{k\ge0}(-1)^ne_nh_k\partial /\partial h_{n+k}
\\
\phantom{\sum_{n\ge1}(-1)^ne_np_n^{\perp}}
=\sum_{n\ge1}\sum_{0\le k<n}(-1)^{n-k}e_{n-k}
h_k\partial /\partial h_n=-\sum_{n\ge1}h_n\partial /\partial h_n,
\end{gather*}
where we used the relation~\eqref{pperp} and then~\eqref{eh}.
Therefore by applying the conjugate operator to the vector $h_\lambda$ with $\ell(\lambda)=k$ we get
$- kh_\lambda$.
Now~\eqref{pro1} follows from the duality relation~\eqref{mh}.
\end{proof}

\section{Proof of Theorem~\ref{Theorem3}}\label{Section8}

Denote by $\Phi^{(k)}_n$ the right hand side of the relation~\eqref{psink}.
Then introduce the~series
\begin{gather}
\label{phin}
\Phi_n(u)=\Phi^{(1)}_n/(u+1)_1+\Phi^{(2)}_n/
(u+1)_2+\cdots
\end{gather}
and the column vector
\begin{gather*}
\Phi(u)=
\begin{bmatrix}
\raisebox{0pt}[12pt][0pt]{$\Phi_1(u)$}
\\
\Phi_2(u)
\\
\vdots
\end{bmatrix}
.
\end{gather*}
We have to prove the relation $\Psi(u)=\Phi(u)$.
Multiplying by $u-L$ in the def\/inition~\eqref{psiu} of $\Psi(u)$ the latter relation becomes
\begin{gather*}
p^{\ast}A(u+1)=(u-L)\Phi(u).
\end{gather*}
By taking the entries of the column vectors here and using the def\/inition~\eqref{Ljk} we have to prove
that for every index $n\ge1$
\begin{gather}
\label{pjA}
p^{\ast}_nA(u+1)=(u+n-\alpha n)\Phi_n(u)-\sum_{i\ge1}
p_{i-n}\Phi_i(u).
\end{gather}

By using for $n\ge1$ the expansions~\eqref{au} and~\eqref{phin} along with the obvious relation for $k\ge1$
\begin{gather*}
u/(u+1)_k=1/(u+1)_{k-1}-k/(u+1)_k,
\end{gather*}
the relation~\eqref{pjA} is equivalent to the collection of relations for all $k\ge1$
\begin{gather}
\label{phink}
p^{\ast}_nA^{(k)}=\Phi^{(k+1)}_n+(
n-\alpha n-k)\Phi^{(k)}_n-\sum_{i\ge1}
p_{i-n}\Phi^{(k)}_n
\end{gather}
together with the relation $p^{\ast}_n=\Phi^{(1)}_n$.
But the latter relation follows from the def\/inition of $\Phi^{(1)}_n$ by
using~\eqref{firstfactor}.
Indeed, if $\lambda=\mu\sqcup\nu$ where $\ell(\lambda)=1$ and $|\nu|=n\ge1$, then $\lambda=\nu=(n)$ while
$\mu$ has no parts dif\/ferent from zero.
We then observe that $p_n^{\ast}=m_{(n)}^{\ast}$.
So we have to prove~\eqref{phink}.

For any two symmetric functions $f,g\in\Lambda$ the operator $f g^{\ast}$
acting on the vector space $\Lambda$ can be represented by its \emph{symbol} $f\otimes
g\in\Lambda\otimes\Lambda$.
We will extend this representation to inf\/inite sums of operators of the form $f
g^{\ast}$ by linearity.
In particular, the symbol of the operator at the right hand side of~\eqref{basic} is the inf\/inite linear
combination of elements of the vector space $\Lambda\otimes\Lambda$
\begin{gather}
\label{dek}
\Delta^{(k)}=(-1)^k\sum_{\ell(\lambda)=k}d_\lambda m_\lambda\otimes m_\lambda.
\end{gather}

Extending our notation, set $\Phi^{(k)}_0=A^{(k)}$.
By Theorem~\ref{Theorem1} this is indeed the right hand side of~\eqref{psink} at $n=0$.
Thus for any $n\ge0$ the symbol of the operator $\Phi^{(k)}_n$ equals $(-1)^{
n}(e_n^{\perp}\otimes1)\Delta^{(k)}$.

Let us now prove~\eqref{phink} for $n,k\ge1$.
Using our extended notation,~\eqref{phink} can be rewritten~as
\begin{gather*}
\sum_{1\le i\le n}p^{\ast}_i\Phi_{n-i}^{(k)}+\sum_{i\ge1}
p_i\Phi_{n+i}^{(k)}=\Phi^{(k+1)}_n+(
n-\alpha n-k)\Phi_n^{(k)}.
\end{gather*}
In the representation by symbols the latter operator relation becomes
\begin{gather}
\sum_{1\le i\le n}(-1)^{n-i}\big(\alpha p^{\perp}_i e_{n-i}^{\perp}\otimes1+e_{n-i}^{\perp}\otimes p_i\big)\Delta^{(k)}
+\sum_{i\ge1}(-1)^{n+i}\big(p_ie_{n+i}^{\perp}\otimes1\big)\Delta^{(k)}
\nonumber
\\
\qquad {}
=(-1)^{n}\big(e_n^{\perp}\otimes1\big)\Delta^{(k+1)}+(n-\alpha n-k)(-1)^{n}\big(e_n^{\perp}\otimes1\big)\Delta^{(k)}.
\label{symb}
\end{gather}
Here we take commutators with $p^{\ast}_i$ and then use the equality
$p^{\ast}_i=\alpha p^{\perp}_i$ given by~\eqref{pnast}.

Observe that both sides of the relation~\eqref{symb} are polynomials in $\alpha$ of degree
$1$.
By equating their coef\/f\/icients at~$\alpha$ we get the relation
\begin{gather*}
\sum_{1\le i\le n}(-1)^{n-i}\big(p^{\perp}_i e_{n-i}
^{\perp}\otimes1\big)\Delta^{(k)}=n(-1)^{n+1}
\big(e_n^{\perp}\otimes1\big)\Delta^{(k)}.
\end{gather*}
But this relation follows from the identity
\begin{gather*}
\sum_{1\le i\le n}(-1)^{n-i}e_{n-i}p_i=n(-1)^{n+1}e_n,
\end{gather*}
which holds by~\eqref{hep}.
It now remains to verify the relation~\eqref{symb} only for $\alpha=0$, that is
\begin{gather}
\sum_{1\le i\le n}(-1)^{n-i}\big(e_{n-i}^{\perp}
\otimes p_i\big)\Delta^{(k)}+\sum_{i\ge1}(-1)^{n+i}
\big(p_ie_{n+i}^{\perp}\otimes1\big)\Delta^{(k)}
\nonumber
\\
\qquad{}
=(-1)^{n}\big(e_n^{\perp}\otimes1\big)\Delta^{(k+1)}+(
n-k)(-1)^{n}\big(e_n^{\perp}\otimes1\big)\Delta^{(k)}.
\label{symbzero}
\end{gather}

It follows from the def\/initions~\eqref{mon} and~\eqref{pon} that for any integer $i\ge1$ and any
partition $\mu$
\begin{gather}\label{pim}
p_im_\mu=\sum_\lambda c_{\lambda\mu}m_\lambda,
\end{gather}
where $\lambda$ ranges over all those partitions which can be obtained by increasing by $i$ anyone of the
\emph{distinct} parts of $\mu$.
That includes increasing a~zero part of $\mu$, in which case $\ell(\lambda)=\ell(\mu)+1$.
Otherwise we have $\ell(\lambda)=\ell(\mu)$ of course.
In any case, the coef\/f\/icient $c_{\lambda\mu}$ in~\eqref{pim} is equal to the
multiplicity in $\lambda$ of that part which has been obtained by increasing.
Note that if the increased part of $\mu$ is zero, that is if $\lambda=\mu\sqcup i$, then
the product $c_{\lambda\mu}d_\mu$ equals $d_\lambda$.
But if the increased part of $\mu$ is not zero, then the ratio $
c_{\lambda\mu}d_\mu/ d_\lambda $ equals the multiplicity of
the increased part of $\mu$.
Let us denote the latter multiplicity by $c_{\mu\lambda}$, so that we have
$c_{\lambda\mu}d_\mu=c_{\mu\lambda}d_\lambda$ when
$\ell(\lambda)=\ell(\mu)$.

The expansion~\eqref{pim} also shows that either side of the relation~\eqref{symbzero} can be written as
a~sum of tensor products of the form $f\otimes m_\lambda$ where $f\in\Lambda$ and $\ell(\lambda)=k,k+1$;
see the def\/inition~\eqref{dek}.
By taking only the summands $f\otimes m_\lambda$ in~\eqref{symbzero} where
$\ell(\lambda)=k+1$, we get the relation
\begin{gather*}
\sum_{1\le i\le n}(-1)^{n-i+k}\sum_{\ell(\mu)=k}c_{\mu
\sqcup i,\mu}d_\mu e_{n-i}^{\perp}(m_\mu)\otimes m_{\mu\sqcup i}
=(-1)^{n+k+1}\sum_{\ell(\lambda)=k+1}d_\lambda e_n^{\perp}(m_\lambda)\otimes m_\lambda.
\end{gather*}
By using~\eqref{hperp} and the equality $ c_{\mu\sqcup
i,\mu}d_\mu=d_{\mu\sqcup i}
$ the last displayed relation be rewritten as
\begin{gather*}
\sum_{1\le i\le n}(-1)^{n-i+k}\sum_{\ell(\lambda)=k+1}d_\lambda e_{n-i}
^{\perp}h_i^{\perp}(
m_\lambda)\otimes m_\lambda=(-1)^{n+k+1}\sum_{\ell(\lambda)=k+1}
d_\lambda e_n^{\perp}(m_\lambda)\otimes m_\lambda.
\end{gather*}
But this relation immediately follows from an identity which holds by~\eqref{eh},
\begin{gather*}
\sum_{1\le i\le n}(-1)^{n-i}h_ie_{n-i}=(-1)^{n+1}e_n.
\end{gather*}

Now take the summands $f\otimes m_\lambda$ in~\eqref{symbzero} where $\ell(\lambda)=k$.
In this way we get the relation
\begin{gather}
\sum_{1\le i\le n}(-1)^{n-i+k}\sum_{\ell(\mu)=k}\sum_\lambda
\qquad
c_{\lambda\mu}d_\mu e_{n-i}^{\perp}(m_\mu)\otimes m_\lambda
\nonumber
\\
\qquad\phantom{=}
{}+\sum_{i\ge1}\sum_{\ell(\lambda)=k}(-1)^{n+i+k}d_\lambda
p_ie_{n+i}^{\perp}(m_\lambda)\otimes m_\lambda
\nonumber
\\
\qquad{}
=(n-k)(-1)^{n+k}\sum_{\ell(\lambda)=k}d_\lambda e_n^{\perp}(m_\lambda)\otimes m_\lambda,
\label{threeline}
\end{gather}
where in the f\/irst of the three displayed lines the partition $\lambda$ is obtained by increasing by $i$
any of the distinct non-zero parts of $\mu$.
Let us consider the sum in the second line of the display~\eqref{threeline}.
\begin{proposition}\label{Prop2}
For any $n\ge1$ we have the operator identity
\begin{gather}
\label{prop2}
\sum_{i\ge1}(-1)^{n+i}p_ie_{n+i}^{\perp}
=\sum_{0\le i\le n}\sum_{j\ge1}(-1)^{n-i+1}(
\partial /\partial h_j)^{\perp}e_{n-i}
^{\perp}h_{i+j}^{\perp}.
\end{gather}
\end{proposition}
\begin{proof}
Using the form $\langle\;,\;\rangle$ let us conjugate the operator at the left hand side of~\eqref{prop2}.
We get
\begin{gather*}
\sum_{i\ge1}(-1)^{n+i}e_{n+i}p_i^{\perp}
=\sum_{i\ge1}\sum_{l\ge0}(-1)^{n+i}e_{n+i}h_{l}\partial /\partial h_{i+l}
=\sum_{j\ge1}\sum_{0\le l<j}(-1)^{n+j-l}e_{n+j-l}h_{l}\partial /\partial h_{j}
\\
\phantom{\sum_{i\ge1}(-1)^{n+i}e_{n+i}p_i^{\perp}}
=\sum_{j\ge1}\sum_{j\le l\le n+j}(-1)^{n+j-l+1}e_{n+j-l}h_{l}\partial /\partial h_{j}
\\
\phantom{\sum_{i\ge1}(-1)^{n+i}e_{n+i}p_i^{\perp}}
=\sum_{j\ge1}\sum_{0\le i\le n}(-1)^{n-i+1}e_{n-i}h_{i+j}\partial /\partial h_{j}
\end{gather*}
as required.
Here we use~\eqref{pperp}, then denote $j=i+l$, then use~\eqref{eh}, and f\/inally set
$i=l-j$.
\end{proof}

Further, the right hand side of the identity~\eqref{prop2} can be rewritten as
\begin{gather*}
\sum_{0\le i\le n}\sum_{j\ge1}(-1)^{n-i+1}e_{n-i}^{\perp}(\partial /\partial h_j)^{\perp}
h_{i+j}^{\perp}+\sum_{0\le i\le n}\sum_{j\ge1}(-1)^{n-i}(\partial e_{n-i}/\partial h_j)^{\perp}h_{i+j}^{\perp},
\end{gather*}
where we use the commutator of the operator $(\partial /\partial
h_j)^{\perp}$ with $e_{ n-i}^{\perp}$.
Consider the conjugate of the second group of summands in the above display.
The conjugate is an operator of multiplication by the element of $\Lambda$ which is computed in the next
proposition.
\begin{proposition}\label{Prop3}
For any $n\ge1$ we have the identity in $\Lambda$
\begin{gather}
\label{prop3}
\sum_{0\le i\le n}\sum_{j\ge1}(-1)^{n-i}h_{i+j}
\partial e_{n-i}/\partial
h_j=(-1)^nne_n.
\end{gather}
\end{proposition}
\begin{proof}
The summation at the left hand side of~\eqref{prop3} can be restricted to $0\le i<n$~without af\/fecting
the value of the sum, because $e_0=1$.
Then by using~\eqref{ehnu} the sum can be rewritten as
\begin{gather}
\sum_{0\le i<n}\sum_{j\ge1}h_{i+j}\sum_{|\mu|=n-i}(-1)^{\ell(\mu)}
\partial h_\mu/\partial h_j\cdot\ell(\mu)
!/d_\mu
\nonumber
\\
\qquad\qquad
=\sum_{0\le i<n}\sum_{j\ge1}\sum_{|\mu|=n-i}(-1)^{\ell(\mu)}h_\nu\cdot c_{\mu\nu}\ell(\mu)!/d_\mu,
\label{prop4}
\end{gather}
where $\mu$ ranges over all partitions of $n-i$ containing $j$ as a~part.
Further, here $\nu$ denotes the partition obtained from $\mu$ by replacing a~part $j$ by $i+j$.
According to the notation introduced after stating~\eqref{pim}, here $c_{\mu\nu}$ is the
multiplicity of the part $j$ in $\mu$.
In particular, $\ell(\mu)=\ell(\nu)$ and
\begin{gather}
\label{cd}
c_{\mu\nu}/d_\mu=c_{\nu\mu}
/d_\nu.
\end{gather}
Here \looseness=1 we extend the notation $c_{\mu\nu}$ to the case $i=0$.
Then $\mu=\nu$ and the part $j$ is not determined by the notation.
This should not cause us any confusion however, as it will always remain clear which multiplicity is taken.
In particular, the equality~\eqref{cd} will remain valid in the case $i=0$.

By again using~\eqref{ehnu} the sum at the right hand side of~\eqref{prop3} can be rewritten as
\begin{gather}
\label{prop5}
\sum_{|\nu|=n}(-1)^{\ell(\nu)}h_\nu\cdot n\ell(\nu)!/d_\nu
\end{gather}
where \looseness=1 $\nu$ ranges over all partitions of $n$.
To prove the identity~\eqref{prop3} it remains to observe that for any f\/ixed $\nu$ the terms at the right
hand side of~\eqref{prop4} containing the factor $h_\nu$ can be collected as follows.
Choose any of the distinct non-zero parts of $\nu$ and replace it by a~positive integer not exceeding the
original part, hence getting some $\mu$.
For any choice of the original part there are as many possible replacements as the size of the part.
By multiplying the size by the ratio~\eqref{cd} and then summing over all distinct non-zero parts of $\nu$
we get $n/d_\nu$, exactly as in~\eqref{prop5}.
\end{proof}

We can now complete our proof of Theorem~\ref{Theorem3}.
By using Proposition~\ref{Prop2} and then Proposition~\ref{Prop3}
as explained just before stating the latter, the sum in the
second line of the display~\eqref{threeline} equals
\begin{gather}
\sum_{0\le i\le n}\sum_{j\ge1}\sum_{\ell(\lambda)=k}(-1)^{n-i+k+1}
d_\lambda e_{n-i}^{\perp}(\partial /\partial h_j)^{\perp}h_{i+j}^{\perp}(m_\lambda)\otimes m_\lambda
\nonumber
\\
\qquad\qquad
{}+\sum_{\ell(\lambda)=k}(-1)^{n+k}ne_n^{\perp}(m_\lambda)\otimes m_\lambda.
\label{last}
\end{gather}
In \looseness=1 the f\/irst line of the above display the summation can be restricted to partitions $\lambda$ containing
$i+j$ as a~part, without af\/fecting the value of the sum; see~\eqref{hperp}.
For these partitions $\lambda$ we have
\begin{gather*}
d_\lambda(\partial /\partial h_j)^{\perp}
h_{i+j}^{\perp}(m_\lambda)=d_\lambda c_{\mu\lambda}m_\mu=d_\mu c_{\lambda\mu}m_\mu,
\end{gather*}
where $\mu$ denotes the partition obtained from $\lambda$ by replacing a~part $i+j$ with $j$; see the
remark we made immediately after displaying~\eqref{cd}.
Therefore the summands in~\eqref{last} with $1\le i\le n$ cancel with the sum in the f\/irst line
of~\eqref{threeline}.
Since the sum of the multiplicities of non-zero parts of $\lambda$ is $\ell(\lambda)$, the summands
in~\eqref{last} with $i=0$ add up to
\begin{gather*}
\sum_{\ell(\lambda)=k}(-1)^{n+k+1}d_\lambda ke_n^{\perp}(m_\lambda)\otimes m_\lambda.
\end{gather*}
The \looseness=1 latter sum and the sum in the second line of~\eqref{last} cancel with the right hand side
of~\eqref{threeline}.
All these cancellations establish the relation~\eqref{threeline}.
Our proof of Theorem~\ref{Theorem3} is now completed.

